\clearpage
\thispagestyle{empty}\SourceLeaf{283}
\pdfbookmark[0]{Anzeigen eigener Abhandlungen}{anzeigen-title}
\begin{center}
\vspace*{2in}
{\Large ANZEIGEN\par}
\vspace{0.5in}
EIGENER\par
\vspace{0.5in}
{\Large ABHANDLUNGEN.\par}
\end{center}
\clearpage
\SourceLeaf{285}
\fancyhead[L]{\small Anzeigen eigener Abhandlungen}
\pdfbookmark[1]{Göttingische gelehrte Anzeigen. 5. April 1813.}{anzeige-1813-april5}
\begin{center}
\rule{0.85\textwidth}{0.3pt}\par
Göttingische gelehrte Anzeigen. 5. April 1813.\par
\rule{0.85\textwidth}{0.3pt}
\end{center}
\medskip
Am 18. März reichte Professor \textsc{Gauss} bei der Königlichen Sozietät einen Vortrag mit dem folgenden Titel ein:
\begin{center}
\textit{Theoria attractionis corporum sphaeroidicorum ellipticorum homogeneorum, methodo nova tractata.}
\end{center}
Bekanntlich haben sich seit \textsc{Newton}s Zeiten die bedeutendsten Geometer mit dieser berühmten Aufgabe beschäftigt, beinahe im Wettstreit miteinander. \textsc{Newton} tat den ersten Schritt, indem er zeigte, wie sich die Anziehung auf einen Punkt auf der Achse eines Rotationssphäroids bestimmen lässt. Außerdem zeigte er, in welchem einfachen Verhältnis die Anziehungen auf alle Punkte stehen, die im Innern des Sphäroids auf ein und demselben Durchmesser liegen. Danach gelang es \textsc{Maclaurin}, in seiner berühmten Preisschrift über Ebbe und Flut die Anziehung auf alle Punkte der Oberfläche des Sphäroids zu bestimmen. Nach \textsc{Newton}s Lehrsatz ließen sich auch die Anziehungen im Innern darauf zurückführen, sodass nur noch die Anziehung auf die äußeren Punkte zu bestimmen war. Das war freilich der schwierigste Teil der Aufgabe. Auch hier tat \textsc{Maclaurin} bereits einen Schritt. Er bestimmte die Anziehung auf die Punkte in der Verlängerung des Äquators und der Achse. \textsc{Maclaurin}s Entdeckungen wurden allgemein als Meisterwerke der Synthese bewundert und eine Zeit lang als Beweise dafür angesehen, dass es Fälle gibt, in denen die synthetische Methode einen entschiedenen Vorzug vor
\SourceLeaf{286}der analytischen hat. \textsc{Lagrange} verhalf der analytischen Methode wieder zu ihrem Recht, indem er mit ihr eine Aufgabe löste, die nur der Synthese zugänglich erschienen war, und mit der ihm eigenen Gewandtheit zeigte, wie sich alle Entdeckungen \textsc{Maclaurin}s auf analytischem Weg gewinnen lassen. Obwohl dies in der Sache selbst keinen neuen Fortschritt brachte, musste es doch als höchst wichtige Vorbereitung für die späteren Arbeiten gelten. \textsc{Legendre} gelang es, die Theorie der Anziehung von Rotationssphäroiden zu vollenden, indem er den schönen Lehrsatz fand und bewies. Haben die beiden erzeugenden Ellipsen denselben Brennpunkt, so hat die Anziehung eines Sphäroids auf einen äußeren Punkt dieselbe Richtung wie die Anziehung auf denselben Punkt durch ein zweites Sphäroid, dessen Oberfläche durch diesen Punkt geht. Die erste Anziehung verhält sich zur zweiten wie die Masse des ersten Sphäroids zur Masse des zweiten.

All dies bezieht sich auf Sphäroide, die durch Drehung einer halben Ellipse um die eine oder die andere Achse entstehen. Nun blieb aber noch die weit schwierigere Aufgabe, die Anziehung eines Ellipsoids zu bestimmen, bei dem auch der Äquator elliptisch ist, also eines Körpers, dessen jeder Schnitt mit einer Ebene eine Ellipse ergibt. Die Bestimmung der Anziehung auf Punkte in Richtung der drei Hauptachsen hatte bereits \textsc{Maclaurin} angedeutet, und \textsc{d'Alembert} und \textsc{Lagrange} hatten dafür analytische Beweise gegeben. Außerdem hatte \textsc{Legendre} durch Induktion die allgemeine Gültigkeit seines eben angeführten schönen Lehrsatzes vermutet, ohne einen strengen Beweis finden zu können. Es blieb \textsc{Laplace} vorbehalten, diese Lücke zu füllen und die Lösung der Aufgabe in ihrer ganzen Allgemeinheit zu vollenden (1782). Damit, könnte man glauben, war die Untersuchung abgeschlossen. Doch schon der Umstand, dass sich seitdem mehrere Geometer erneut mit demselben Gegenstand beschäftigt haben, zeigt, dass noch viel zu wünschen übrig blieb. Das Erste und Wesentlichste bei einer Aufgabe ist immer, dass sie überhaupt gelöst wird. Zum selben Ziel führen aber oft mehrere Wege. Nicht selten gelangt man zunächst auf einem langen, dornigen Umweg ans Ziel. Der kürzeste, wahre, echte Weg wird erst viel später entdeckt. \textsc{Laplace}s Lösung ist ein schönes Zeugnis feinster analytischer Kunst. Doch der Weg, auf dem er zu ihr gelangt. ist lang und beschwerlich, und gewiss sind ihm nur wenige Geometer und
\SourceLeaf{287}Astronomen auf diesem Weg gefolgt. Auch der Gebrauch unendlicher Reihen, deren Konvergenz nicht bewiesen ist, beeinträchtigt die Klarheit und Schlüssigkeit des Beweises ein wenig. \textsc{Legendre} hat zwar 1788 eine andere Lösung gegeben, für die jedoch fast dasselbe gilt, was wir gegen die Lösung \textsc{Laplace}s eingewandt haben. Ein sachkundiger Richter, \textsc{Lagrange}, fällt über die Lösungen jener beiden großen Analytiker das folgende Urteil (in den \textit{Nouv. Mém. de l'Acad. de Berlin} 1793): ‘\textit{On ne peut regarder leurs solutions que comme des chefs-d'oeuvre d'analyse, mais on peut désirer encore une solution plus directe et plus simple: et les progrès continuels de l'analyse donnent lieu de l'espérer.}’ Seitdem haben auch \textsc{Biot} und \textsc{Plana} versucht, jene beiden Beweise zu verbessern und zu vereinfachen. So wertvoll diese Arbeiten sind, muss man diese Lösungen doch noch immer zu den verwickeltsten und subtilsten Anwendungen der Analysis rechnen.

Der Verfasser der vorliegenden Abhandlung war seit Langem davon überzeugt, dass die \textit{echte} Methode zur Lösung jener berühmten Aufgabe erst noch gefunden werden musste. Vor einem halben Jahr erhielt er Anlass, sich näher mit ihr zu beschäftigen. Er schlug einen von den bisherigen ganz verschiedenen Weg ein und hatte das Vergnügen, eine so überraschend kurze und einfache Lösung zu finden, dass sich das Wesentliche auf zwei Seiten darstellen ließ. Freilich hat er sie hier etwas ausführlicher vorgetragen. Zum einen wollte er sie auch weniger geübten Lesern verständlich machen, denen diese für die Gestalt der Erde so interessanten Untersuchungen bisher völlig unzugänglich waren. Dafür, dass sich die neue Lösung dazu uneingeschränkt eignet, hat er bereits mehrere Beweise. Zum anderen schien es der Mühe wert, die Gründe, auf denen sie beruht und die sich auch bei anderen Gelegenheiten oft mit Vorteil anwenden lassen, etwas ausführlicher darzulegen, als es für den unmittelbaren Zweck erforderlich gewesen wäre.

Wir wollen nun die Hauptpunkte der ganzen Lösung so kurz wie möglich darstellen, dabei aber so vollständig, wie es für Kenner nötig ist. Wir müssen hier darauf verzichten, auch für Leser, die mit Untersuchungen dieser Art noch nicht vertraut sind, ganz verständlich zu sein. Diese müssen wir auf die ausführliche Abhandlung selbst verweisen, die bereits gedruckt ist und in Kürze im zweiten Band der \textit{Commentationes recentiores} der Sozietät erscheinen wird.

\SourceLeaf{288}Der Verfasser beginnt damit, sechs verschiedene allgemeine Lehrsätze zu begründen, mit deren Hilfe sich dreifache Integrale über einen räumlichen Bereich auf zweifache Integrale allein über die Oberfläche des Körpers zurückführen lassen. Wir geben hier nur drei dieser Lehrsätze an, da die anderen für die vorliegende Untersuchung nicht nötig sind.

Es sei $\mathrm{d}s$ ein Element der Oberfläche eines Körpers von beliebiger Gestalt. Es seien $PQ,PM,PX,PY,PZ$ die Geraden, die von einem Punkt $P$ dieses Elements senkrecht zur Oberfläche nach außen, zum angezogenen Punkt $M$ und parallel zu den drei Koordinatenachsen gezogen werden. Ferner sei $r$ der Abstand des Punktes $M$ von $P$, $MQ$ der Winkel zwischen $PM$ und $PQ$, $MX$ der Winkel zwischen $PM$ und $PX$ und $QX$ der Winkel zwischen $PQ$ und $PX$. Schließlich bezeichne $\pi$ das Verhältnis des Kreisumfangs zum Durchmesser und $X$ die Anziehung, die der ganze Körper auf den Punkt $M$ parallel zur Koordinatenachse $x$ ausübt. Dann gilt
\[
\text{I.}\qquad \int\frac{\mathrm{d}s.\cos MQ}{rr}=0\text{ oder }=-4\pi
\]
je nachdem, ob $M$ außerhalb oder innerhalb des Körpers liegt, und
\[
\begin{aligned}
\text{II.}\qquad &\int\frac{\mathrm{d}s.\cos QX}{r}=X\\
\text{III.}\qquad &\int\frac{\mathrm{d}s.\cos MQ.\cos MX}{r}=-X
\end{aligned}
\]
wobei die Integrale über die ganze Oberfläche des Körpers zu erstrecken sind. Wir übergehen hier die Beweise dieser Lehrsätze und bemerken nur, dass die ersten beiden auf der Zerlegung des Körpers in Kegelelemente mit der Spitze in $M$ beruhen, der dritte hingegen auf der Zerlegung in prismatische Elemente parallel zur Koordinatenachse $x$.

Für die Oberfläche eines Ellipsoids mit den drei Halbachsen $A,B,C$ gilt zwischen den Koordinaten $x,y,z$ die Gleichung
\[
\frac{xx}{AA}+\frac{yy}{BB}+\frac{zz}{CC}=1
\]
Ferner gilt $\cos QX=\frac{x}{AA\rho}$, wenn man der Kürze halber setzt
\[
\BareRoot{(\frac{xx}{A^4}+\frac{yy}{B^4}+\frac{zz}{C^4})}=\rho
\]
Bezeichnen $a,b,c$ die Koordinaten des Punktes $M$, so gilt
\SourceLeaf{289}
\[
\begin{aligned}
r&=\BareRoot{[(a-x)^2+(b-y)^2+(c-z)^2]}\\
\cos MX&=\frac{a-x}{r}\\
\cos MQ&=\frac{1}{\rho r}(\frac{(a-x)x}{AA}+\frac{(b-y)y}{BB}+\frac{(c-z)z}{CC})
\end{aligned}
\]
Nun werden zwei neue Veränderliche $p,q$ eingeführt, von denen $x,y,z$ folgendermaßen abhängen.
\[
\begin{aligned}
x&=A\cos p\\
y&=B\sin p\cos q\\
z&=C\sin p\sin q
\end{aligned}
\]
Um also die ganze Oberfläche des Ellipsoids zu erfassen, muss man $p$ von $0$ bis $180^0$ und $q$ von $0$ bis $360^0$ laufen lassen. Schließlich setze man noch $X=ABC\xi$. Aus bekannten Gründen ergibt sich $\mathrm{d}s=\mathrm{d}p.\mathrm{d}q.ABC\rho\sin p$. Die drei obigen Lehrsätze nehmen damit die folgende Gestalt an, wenn man der Kürze halber
\[
\frac{(a-x)x}{AA}+\frac{(b-y)y}{BB}+\frac{(c-z)z}{CC}=\psi
\]
setzt.
\[
\begin{aligned}
[1]\qquad &\int\int\frac{\mathrm{d}p.\mathrm{d}q.\sin p.\psi}{r^3}=0\text{ oder }=-\frac{4\pi}{ABC}\\
[2]\qquad &\int\int\frac{\mathrm{d}p.\mathrm{d}q.\cos p.\sin p}{r}=A\xi\\
[3]\qquad &\int\int\frac{\mathrm{d}p.\mathrm{d}q.\sin p.\psi(a-x)}{r^3}=-\xi
\end{aligned}
\]
Man betrachte nun $A,B,C$ als bestimmte einzelne Werte dreier Veränderlicher $\alpha,\beta,\gamma$, die jedoch so miteinander verbunden sind, dass $\alpha\alpha-\beta\beta$ und $\alpha\alpha-\gamma\gamma$ konstant bleiben. Aus Formel [1] folgt leicht, dass bei unbegrenzt wachsendem $\alpha$ die Größe $\xi$ unbegrenzt abnimmt. Differenziert man [2] nach den Veränderlichen $\alpha,\beta,\gamma$ und verwendet dabei das Variationszeichen $\delta$, so erhält man
\[
\begin{aligned}
\alpha\delta\xi+\xi\delta\alpha&=-\int\int\frac{\mathrm{d}p.\mathrm{d}q.\cos p.\sin p.\delta r}{rr}\\
&=\delta\alpha.\int\int\frac{\mathrm{d}p.\mathrm{d}q.\sin p.x\psi}{r^3}
\end{aligned}
\]
\SourceLeaf{290}oder, wenn man hier für $\xi$ seinen Wert aus [3] einsetzt,
\[
\alpha\delta\xi=\delta\alpha.\int\int\frac{\mathrm{d}p.\mathrm{d}q.\sin p.a\psi}{r^3}
\]

Der Vergleich mit [1] ergibt
\[
[4]\qquad\delta\xi=0
\]
wenn der Punkt $M$ außerhalb des Ellipsoids liegt, und
\[
[5]\qquad\delta\xi=-\frac{4\pi a\delta\alpha}{\alpha\alpha\beta\gamma}
\]
wenn $M$ innerhalb liegt.

Aus [4] folgt, dass $\xi$ konstant ist oder dass die Anziehung der Masse proportional ist, und zwar für alle Sphäroide, deren Hauptschnitte Ellipsen mit denselben Brennpunkten sind, solange $M$ nicht innerhalb liegt. Die Bestimmung der Anziehung eines Sphäroids auf einen äußeren Punkt führt also auf die Bestimmung der Anziehung eines anderen Sphäroids zurück, das aus denselben Brennpunkten beschrieben wird und durch den angezogenen Punkt geht. Um diese Anziehung zu bestimmen, betrachte man den anderen Fall, in dem der angezogene Punkt innerhalb liegt. Setzt man
\[
\begin{aligned}
\beta\beta&=\alpha\alpha-AA+BB\\
\gamma\gamma&=\alpha\alpha-AA+CC
\end{aligned}
\]
in Gleichung [5] ein, setzt zugleich $\frac{A}{\alpha}=t$ und verwendet anstelle des Zeichens $\delta$ wieder das gewöhnliche $d$, so erhält man
\[
\xi=\frac{4\pi a}{A^3}\int\frac{tt\mathrm{d}t}{\BareRoot{[(1-(1-\frac{BB}{AA})tt)(1-(1-\frac{CC}{AA})tt)]}}
\]
wobei das Integral so festgelegt werden muss, dass es für $t=0$ verschwindet und dann für das \textit{bestimmte} Sphäroid bis $t=1$ zu erstrecken ist. Man erhält also im selben Sinne
\[
[6]\qquad X=\frac{4\pi aBC}{AA}\int\frac{tt\mathrm{d}t}{\BareRoot{[(1-(1-\frac{BB}{AA})tt)(1-(1-\frac{CC}{AA})tt)]}}
\]
Diese Formel gibt die Anziehung für alle Punkte an, die nicht außerhalb liegen. Da sie bis zur Oberfläche selbst gelten muss und die Anziehung äuße\SourceLeaf{291}rer Punkte bereits auf die Anziehung der Punkte auf der Oberfläche zurückgeführt wurde, ist die Aufgabe damit vollständig gelöst. (Es braucht kaum erwähnt zu werden, dass sich die Anziehung parallel zu den beiden anderen Hauptachsen jeweils einfach durch Vertauschen von $A,a$ mit $B,b$ oder mit $C,c$ ergibt.)

Gleichung [6] zeigt ferner, dass für einen inneren Punkt die Anziehung aller Sphäroide, die einander ähnlich sind und ähnlich liegen, dieselbe ist. Denkt man sich also ein solches Sphäroid in Schichten unterteilt, die durch ähnliche ellipsoidische Flächen begrenzt sind, so ist klar, dass alle außerhalb des angezogenen Punktes liegenden Schichten überhaupt nichts zur Anziehung beitragen und nur die Anziehung des sphäroidischen Kerns übrig bleibt, dessen Oberfläche durch den angezogenen Punkt geht.

Zum Schluss erwähnt der Verfasser noch die neueste Arbeit über denselben Gegenstand von Herrn \textsc{Ivory} in den \textit{Philos. Transact.} von 1809. Herr Graf \textsc{Laplace} machte ihn darauf aufmerksam, und er lernte sie erst kennen, als seine eigene Abhandlung bereits ganz abgeschlossen war. Mit einem sehr glücklichen Einfall hat Herr \textsc{Ivory} die Anziehung auf einen äußeren Punkt auf die Anziehung auf einen inneren Punkt zurückgeführt. Die Art, wie er die Anziehung auf innere Punkte selbst bestimmt, ist zwar voller Scharfsinn und Kunstfertigkeit, beruht aber, ebenso wie \textsc{Laplace}s Lösung für äußere Punkte, zum Teil auf der Betrachtung unendlicher, nicht überall konvergierender Reihen und bleibt weit hinter der Einfachheit zurück, die man sich hätte wünschen können. Daher ist \textsc{Ivory}s Lösung der Aufgabe im Ganzen betrachtet im Grunde nicht viel weniger kunstvoll und verwickelt als die Lösungen von \textsc{Laplace} und \textsc{Legendre}. Im Übrigen beruhen jene Lösung und die hier besprochene auf ganz verschiedenen Grundlagen und haben gar nichts gemeinsam außer dem Gebrauch der zwei Veränderlichen, die oben mit $p,q$ bezeichnet sind.

